\documentclass[a4paper,12pt]{article}

\usepackage[T2A]{fontenc}
\usepackage[utf8]{inputenc}
\usepackage[russian]{babel}
\usepackage{amsmath,amssymb,amsthm,mathtools}

\usepackage{helvet}
% Helvetica не содержит штатных T2A-начертаний в стандартной поставке pdfLaTeX.
% Для кириллицы используем близкий по характеру CM Sans, сохраняя sans-serif стиль.
\DeclareFontFamily{T2A}{phv}{}
\DeclareFontShape{T2A}{phv}{m}{n}{<->ssub * cmss/m/n}{}
\DeclareFontShape{T2A}{phv}{b}{n}{<->ssub * cmss/bx/n}{}
\DeclareFontShape{T2A}{phv}{m}{it}{<->ssub * cmss/m/sl}{}
\DeclareFontShape{T2A}{phv}{b}{it}{<->ssub * cmss/bx/sl}{}
\renewcommand{\familydefault}{\sfdefault}

\usepackage{icomma}
\usepackage[a4paper,margin=20mm,headheight=25pt]{geometry}
\usepackage{microtype}

\usepackage{tikz}
\usetikzlibrary{calc}
\usetikzlibrary{arrows.meta,positioning,shapes.geometric}

\usepackage{tkz-euclide}
\usepackage{pgfplots}
\pgfplotsset{compat=1.18}

\usepackage{tabularx,booktabs,longtable}
\usepackage{enumitem}
\usepackage{hyperref}
\usepackage{parskip}
\usepackage{fancyhdr}
\usepackage{setspace}
\usepackage{xcolor}

% -------------------------------------------------
% Палитра
% -------------------------------------------------

\definecolor{accent}{HTML}{008080}
\definecolor{darkaccent}{HTML}{004D4D}
\definecolor{textgray}{HTML}{333333}
\definecolor{softgray}{HTML}{F3F5F5}

\hypersetup{
  colorlinks=true,
  linkcolor=darkaccent,
  urlcolor=darkaccent
}

\color{textgray}
\onehalfspacing
\setlength{\parindent}{0pt}

\setlist[itemize]{
  leftmargin=1.5em,
  itemsep=0.35em,
  topsep=0.35em
}

\setlist[enumerate]{
  leftmargin=1.7em,
  itemsep=0.55em,
  topsep=0.4em
}

\renewcommand{\arraystretch}{1.25}

% -------------------------------------------------
% Колонтитулы
% -------------------------------------------------

\pagestyle{fancy}
\fancyhf{}

\fancyhead[L]{
  \small\textcolor{darkaccent}{Чек-лист по математике}
}

\fancyhead[R]{
  \small\textcolor{textgray}{6 класс}
}

\fancyfoot[C]{
  \small\thepage
}

\renewcommand{\headrulewidth}{0.4pt}
\renewcommand{\headrule}{
  \hbox to\headwidth{
    \color{accent}
    \leaders\hrule height \headrulewidth
    \hfill
  }
}

% -------------------------------------------------
% Служебные команды
% -------------------------------------------------

\newcommand{\mygcd}[2]{
  \mathop{\text{НОД}}\nolimits(#1,#2)
}

\newcommand{\term}[1]{
  \textbf{\textcolor{darkaccent}{#1}}
}

\newcommand{\exampletitle}[1]{
  \medskip
  \textbf{\textcolor{accent}{#1}}
  \par
  \smallskip
}

% -------------------------------------------------
% Стили блок-схем
% -------------------------------------------------

\tikzset{
  flowstart/.style={
    ellipse,
    draw=accent,
    very thick,
    align=center,
    minimum width=44mm,
    minimum height=10mm,
    fill=softgray
  },
  flowaction/.style={
    rounded corners=3pt,
    rectangle,
    draw=accent,
    very thick,
    align=center,
    text width=112mm,
    minimum height=13mm,
    fill=white
  },
  flowdecision/.style={
    diamond,
    aspect=2.5,
    draw=accent,
    very thick,
    align=center,
    text width=54mm,
    inner sep=2pt,
    fill=softgray
  },
  flowarrow/.style={
    -{Stealth[length=3mm]},
    very thick,
    draw=darkaccent
  },
  flowlabel/.style={
    font=\small,
    fill=white,
    inner sep=1.5pt,
    text=darkaccent
  }
}

\begin{document}

% =================================================
% ТИТУЛЬНЫЙ ЛИСТ
% =================================================

\begin{titlepage}
\thispagestyle{empty}
\centering

\vspace*{28mm}

{\Large
\textcolor{darkaccent}{
\textbf{Чек-лист}
}
\par}

\vspace{8mm}

{\huge\bfseries
Взаимно простые числа, НОД\\
и задачи на одинаковые группы
\par}

\vspace{5mm}

{\Large
\textcolor{accent}{
с коротким введением в факториал
и подсчёт размещений
}
\par}

\vspace{20mm}

{\large 6 класс\par}

\vspace{4mm}

{\large Дата: \today\par}

\vfill

{\large\bfseries
Лёвин Артём Александрович
\par}

\vspace{2mm}

{\large
эксклюзивно для Над Клименко
\par}

\vspace{20mm}

\begin{tikzpicture}[remember picture,overlay]
  \draw[
    accent,
    line width=1.1pt
  ]
    ($(current page.south west)+(18mm,19mm)$)
    ..
    controls
    ($(current page.south west)+(70mm,31mm)$)
    and
    ($(current page.south east)+(-70mm,8mm)$)
    ..
    ($(current page.south east)+(-18mm,20mm)$);
\end{tikzpicture}

\end{titlepage}

% =================================================
% ОСНОВНОЙ ЧЕК-ЛИСТ
% =================================================

\section*{1. Взаимно простые числа}

\term{Определение.}
Два натуральных числа называются взаимно простыми,
если их единственный общий положительный делитель равен~1.

Эквивалентная запись:

\[
\mygcd{a}{b}=1.
\]

\exampletitle{
Пример 1. Проверяем пары среди чисел
$6$, $15$, $30$, $77$
}

\[
\mygcd{6}{77}=1,
\qquad
\mygcd{15}{77}=1,
\qquad
\mygcd{30}{77}=1.
\]

Следовательно, взаимно простые пары:

\[
(6,77),
\qquad
(15,77),
\qquad
(30,77).
\]

Пары
\[
(6,15),\qquad (6,30),\qquad (15,30)
\]
имеют общие делители, большие~1.

% -------------------------------------------------

\subsection*{Правильная дробь и взаимная простота}

Дробь

\[
\frac{m}{n}
\]

является \term{правильной}, если

\[
0<m<n.
\]

Если дополнительно требуется,
чтобы числитель и знаменатель были взаимно простыми,
проверяем:

\[
\mygcd{m}{n}=1.
\]

\exampletitle{
Пример 2. Знаменатель равен $16$
}

Число

\[
16=2^4.
\]

Поэтому с ним взаимно просты все нечётные числа
от~1 до~15:

\[
\frac1{16},
\quad
\frac3{16},
\quad
\frac5{16},
\quad
\frac7{16},
\quad
\frac9{16},
\quad
\frac{11}{16},
\quad
\frac{13}{16},
\quad
\frac{15}{16}.
\]

% =================================================

\section*{2. Как находить НОД через разложение на простые множители}

\term{Шаг 1.}
Разложите каждое число на простые множители,
двигаясь от меньших простых делителей к большим.

\term{Шаг 2.}
Выберите простые множители,
которые входят во все разложения.

\term{Шаг 3.}
Если общий множитель повторяется несколько раз,
берите наименьшее число его повторений
среди всех разложений.

\term{Шаг 4.}
Перемножьте выбранные множители.

\exampletitle{
Пример 3. $126$ красных и $210$ синих бусин
}

Разложим числа:

\[
126=2\cdot3^2\cdot7,
\]

\[
210=2\cdot3\cdot5\cdot7.
\]

Общие множители с минимальными показателями:

\[
2\cdot3\cdot7=42.
\]

Следовательно,

\[
\mygcd{126}{210}=42.
\]

Если все бусины раскладывают
в максимально возможное число одинаковых наборов,
получится

\[
42
\]

набора.

В каждом наборе:

\[
126:42=3
\quad
\text{красные бусины},
\]

\[
210:42=5
\quad
\text{синих бусин}.
\]

% =================================================

\section*{3. Как понять смысл НОД в текстовой задаче}

Перед вычислениями задайте главный вопрос:

\begin{center}
\textbf{
какая величина должна быть общей
для всех частей разбиения?
}
\end{center}

\begin{itemize}

  \item
  Если все предметы делят на
  \term{одинаковые наборы}
  и требуется получить максимально возможное число наборов,
  число наборов должно делить каждое исходное количество.

  Максимальное число наборов равно НОД
  исходных количеств.

  \item
  Если мальчиков и девочек делят
  на отдельные группы
  \term{одинакового размера}
  и требуется получить максимально возможный размер группы,
  размер группы должен делить оба количества.

  Максимальный размер группы равен НОД.

  \item
  После нахождения НОД обязательно подпишите его смысл:

  \begin{center}
  \emph{число наборов},
  \qquad
  \emph{размер группы},
  \qquad
  \emph{число одинаковых частей}.
  \end{center}

\end{itemize}

\exampletitle{
Пример 4. Помидоры и огурцы
}

Есть $84$ помидора и $112$ огурцов.

Все овощи раскладывают
в максимально возможное число одинаковых наборов.

\[
84=2^2\cdot3\cdot7,
\]

\[
112=2^4\cdot7.
\]

Следовательно,

\[
\mygcd{84}{112}
=
2^2\cdot7
=
28.
\]

Получится

\[
28
\]

одинаковых наборов.

В каждом наборе:

\[
84:28=3
\quad
\text{помидора},
\]

\[
112:28=4
\quad
\text{огурца}.
\]

\textbf{
Условие «использовать все предметы»
принципиально важно.
}

Если часть предметов разрешено оставить,
множество допустимых разбиений меняется.

% =================================================
% БЛОК-СХЕМА
% =================================================

\clearpage
\thispagestyle{empty}

\begin{center}
{\Large\bfseries
\textcolor{darkaccent}{
Алгоритм: задача на одинаковые группы или наборы
}
}
\end{center}

\vspace{8mm}

\begin{center}

\begin{tikzpicture}[node distance=7mm]

  \node[flowstart] (start)
  {Начало};

  \node[
    flowaction,
    below=of start
  ] (read)
  {
    Прочитайте условие до конца.
    Выпишите исходные количества
    и все дополнительные ограничения.
  };

  \node[
    flowdecision,
    below=9mm of read
  ] (all)
  {
    Все предметы должны
    быть использованы?
  };

  \node[
    flowaction,
    below=16mm of all,
    xshift=-42mm,
    text width=60mm
  ] (fixed)
  {
    Да:
    исходные количества фиксированы.
    Ищите общий делитель этих чисел.
  };

  \node[
    flowaction,
    below=16mm of all,
    xshift=42mm,
    text width=60mm
  ] (change)
  {
    Разрешено отложить или добавить:
    определите допустимые
    новые количества.
  };

  \node[
    flowaction,
    below=48mm of all
  ] (meaning)
  {
    Определите смысл общего делителя:
    число одинаковых наборов
    или размер одной группы.
  };

  \node[
    flowaction,
    below=of meaning
  ] (gcd)
  {
    Найдите НОД нужных количеств.
    При наличии ограничения на результат
    рассмотрите делители НОД
    и выберите подходящий.
  };

  \node[
    flowaction,
    below=of gcd
  ] (divide)
  {
    Разделите исходные количества
    на выбранное число
    и подпишите смысл полученных частных.
  };

  \node[
    flowstart,
    below=of divide
  ] (finish)
  {
    Проверка условия и ответ
  };

  \draw[flowarrow]
    (start) -- (read);

  \draw[flowarrow]
    (read) -- (all);

  \draw[flowarrow]
    (all)
    --
    node[
      flowlabel,
      pos=0.22,
      left
    ] {да}
    (fixed);

  \draw[flowarrow]
    (all)
    --
    node[
      flowlabel,
      pos=0.22,
      right
    ] {нет}
    (change);

  \draw[flowarrow]
    (fixed) |- (meaning);

  \draw[flowarrow]
    (change) |- (meaning);

  \draw[flowarrow]
    (meaning) -- (gcd);

  \draw[flowarrow]
    (gcd) -- (divide);

  \draw[flowarrow]
    (divide) -- (finish);

\end{tikzpicture}

\end{center}

\clearpage

% =================================================

\section*{4. Дополнительные условия: НОД даёт основу, затем проверяем ограничения}

Иногда найденный НОД
не удовлетворяет дополнительному условию.

В таком случае рассматривают
\term{делители НОД}.

\exampletitle{
Пример 5. Размер группы больше $10$, но меньше $30$
}

Пусть есть $144$ и $192$ человека
двух категорий.

Их распределяют по отдельным группам
одинакового размера.

Сначала:

\[
\mygcd{144}{192}=48.
\]

Все допустимые размеры группы
должны быть делителями числа~$48$:

\[
1,\,
2,\,
3,\,
4,\,
6,\,
8,\,
12,\,
16,\,
24,\,
48.
\]

Дополнительное условие:

\[
10<d<30.
\]

Ему удовлетворяют:

\[
12,\qquad16,\qquad24.
\]

Наибольший допустимый размер группы:

\[
\boxed{24}.
\]

% -------------------------------------------------

\subsection*{Если разрешено отложить часть предметов}

\exampletitle{
Пример 6. $96$ яблок и $145$ груш
}

Разложим числа:

\[
96=2^5\cdot3,
\]

\[
145=5\cdot29.
\]

Поэтому:

\[
\mygcd{96}{145}=1.
\]

При использовании всех предметов
можно составить только один одинаковый набор.

Если разрешено отложить одну грушу,
останется

\[
144
\]

груши.

Тогда:

\[
\mygcd{96}{144}=48.
\]

Можно составить

\[
48
\]

одинаковых наборов.

В каждом наборе:

\[
96:48=2
\quad
\text{яблока},
\]

\[
144:48=3
\quad
\text{груши}.
\]

% =================================================

\section*{5. Многошаговая задача}

\exampletitle{
Пример 7. Наборы распределяют между классами
}

Есть $168$ яблок и $252$ апельсина.

Из всех фруктов составляют
максимально возможное число одинаковых наборов.

\[
168=2^3\cdot3\cdot7,
\]

\[
252=2^2\cdot3^2\cdot7.
\]

Следовательно,

\[
\mygcd{168}{252}
=
2^2\cdot3\cdot7
=
84.
\]

Получили

\[
84
\]

набора.

Их поровну распределяют
между $7$ классами:

\[
84:7=12
\quad
\text{наборов на класс}.
\]

В одном наборе:

\[
168:84=2
\quad
\text{яблока},
\]

\[
252:84=3
\quad
\text{апельсина}.
\]

Поэтому каждый класс получает:

\[
12\cdot2=24
\quad
\text{яблока},
\]

\[
12\cdot3=36
\quad
\text{апельсинов}.
\]

% =================================================

\section*{6. Факториал и размещение людей по местам}

Если $n$ различных людей
рассаживаются на $n$ различных стульев,
то число способов равно:

\[
n\cdot(n-1)\cdot(n-2)\cdots2\cdot1=n!.
\]

Знак

\[
!
\]

называется \term{факториалом}.

\exampletitle{
Пример 8. Девять зрителей и девять стульев
}

На первый стул можно посадить
любого из $9$ зрителей.

После этого на второй стул
остаётся $8$ кандидатов,
затем $7$, $6$ и так далее.

Поэтому:

\[
9!
=
9\cdot8\cdot7\cdot6\cdot5
\cdot4\cdot3\cdot2\cdot1.
\]

Если один заранее определённый зритель
уходит помогать артистам,
остаётся $8$ зрителей и $9$ стульев.

Пустой стул можно выбрать
$9$ способами.

После выбора пустого стула
восьмерых зрителей можно рассадить
по оставшимся местам

\[
8!
\]

способами.

Следовательно:

\[
9\cdot8!=9!.
\]

% =================================================

\section*{7. Типичные ошибки и быстрые проверки}

\begin{itemize}

  \item
  \term{Смысл НОД не подписан.}

  После вычисления обязательно сформулируйте,
  что означает полученное число
  именно в этой задаче.

  \item
  \term{Перепутаны число групп и размер группы.}

  Проверяйте смысл деления:

  \[
  \frac{\text{общее количество}}
       {\text{размер одной группы}}
  =
  \text{число групп},
  \]

  \[
  \frac{\text{общее количество}}
       {\text{число групп}}
  =
  \text{размер одной группы}.
  \]

  \item
  \term{Игнорируется дополнительное ограничение.}

  Условия вида

  \[
  d>10,
  \qquad
  d<30,
  \qquad
  k\ge40
  \]

  требуют дополнительного отбора
  среди делителей
  либо разрешённого изменения исходных данных.

  \item
  \term{Ошибка в разложении на множители.}

  Например,

  \[
  145=5\cdot29.
  \]

  Следовательно, число $145$
  на $3$ не делится.

  \item
  \term{Не проверено использование всех предметов.}

  Если условие допускает остаток,
  возможны другие варианты разбиения.

  \item
  \term{Взаимная простота проверена недостаточно строго.}

  Надёжная проверка:

  \[
  \mygcd{a}{b}=1.
  \]

\end{itemize}

% =================================================
% ТРЕНИРОВОЧНЫЙ БЛОК
% =================================================

\clearpage

\section*{Тренировочный блок --- Путь А}

\textbf{Рекомендация по оформлению.}

Для задач на НОД последовательно подпишите:

\begin{center}
разложения
$\rightarrow$
НОД
$\rightarrow$
смысл НОД
$\rightarrow$
частные
$\rightarrow$
ответ.
\end{center}

\begin{enumerate}


  \item
  Есть $96$ яблок и $145$ груш.

  Нужно составить
  не менее $40$ одинаковых наборов.

  Разрешено отложить несколько груш,
  яблоки требуется использовать все.

  Какое наименьшее число груш
  нужно отложить?

  Сколько наборов получится
  и сколько фруктов каждого вида
  будет в одном наборе?


\end{enumerate}



\end{document}